% !TeX document-id = {a7dcd56d-a58c-47c6-a660-b7c8ed4660db}
% !TeX TXS-program:compile = txs:///pdflatex/[--shell-escape]
% Original Format From https://github.com/koosaga/DeobureoMinkyuParty
% Edited by @evenharder
% These codes should be guaranteed, fast enough, short and easy to type.

\documentclass[landscape, 8pt, a4paper, oneside, twocolumn]{extarticle}

\setlength{\columnseprule}{0.4pt}
\pagenumbering{arabic}

\usepackage{kotex}
\usepackage[hidelinks,unicode,bookmarks=true]{hyperref}
\usepackage[left=2cm,right=2cm,top=2cm,bottom=2cm,a4paper]{geometry}
\usepackage{amsmath}
\usepackage{ulem}
\usepackage{amssymb} 
\usepackage{minted}
\usepackage{color}
\usepackage{titling}
\usepackage{indentfirst}
\usepackage{graphicx}
\usepackage{fancyhdr}
\usepackage{lastpage}
\pagestyle{fancy}

\newcommand{\teamname}{Powered by Zigui}
\lhead{Korea University - \teamname}
\rhead{Page \thepage  \ of \pageref{LastPage} }
\fancyfoot{}

\setlength{\droptitle}{-1cm}

\setminted{breaklines=true, tabsize = 2, breaksymbolleft=}
\usemintedstyle{perldoc}

\title{Team Note of \texttt{\teamname}}
\author{\texttt{@evenharder}(Sangheon Lee), \texttt{@SoulTch}(JEONGJIN LEE), \texttt{@djkim0613}(kim do jae)}
\date{Compiled on \today}

\newcommand{\revised}{Should be \textcolor{red}{\textbf{revised}}.}
\newcommand{\tested}{Should be \textcolor{red}{\textbf{tested}}.}
\newcommand{\added}{Should be \textcolor{red}{\textbf{added}}.}
\newcommand{\WIP}{\textcolor{red}{\textbf{WIP}}}
\begin{document}
{
	\Large
	
	\maketitle
	\setcounter{tocdepth}{2}
    \tableofcontents
}
\thispagestyle{fancy}

\pagebreak

 \textcolor{red}{\textbf{ALL BELOW HERE ARE USELESS IF YOU READ THE STATEMENT WRONG}}
 

\setcounter{section}{-1}
\section{Quotes and Prerequisites}
\begin{verbatim}
evenharder : Mental Abuse To Humans
djkim0613 : 열심히 응원하겠습니다.
SoulTch : How much is this bus ticket?  
* This template is brought from that of 'Deobureo Minkyu Party'
\end{verbatim}
\subsection*{Run script}
\begin{minted}{bash}
#!/bin/bash
g++ -fsanitize=undefined -std=c++14 -O2 -o /tmp/pow $1.cpp
time /tmp/pow < $1.in
# export PATH=~:$PATH
\end{minted}

\subsection*{Debug Code}
\begin{minted}{cpp}
#define setz(x) memset(x, 0, sizeof(x))
#define sz(x) ((int)(x).size())
#define rep(i, e) for (int i = 0, _##i = (e); i < _##i; i++)
#define repp(i, s, e) for (int i = (s), _##i = (e); i < _##i; i++)
#define repr(i, s, e) for (int i = (s)-1, _##i = (e); i >= _##i; i--)
#define repi(i, x) for (auto &i : (x))
// using namespace std;
using ll = long long;
using pii = pair<int, int>;
using pll = pair<ll, ll>;
template<typename T>
ostream &operator<<(ostream &os, const vector<T>& v) {
    cout << "[";
    for (auto p : v) cout << p << ",";
    cout << "]";
    return os;
}
#ifndef __SOULTCH
#define debug(...) 0
#define endl '\n'
#else
#define debug(...) cout << " [-] ", _dbg(#__VA_ARGS__, __VA_ARGS__)
template<class TH> void _dbg(const char *sdbg, TH h){ cout << sdbg << '=' << h << endl; }
template<class TH, class... TA> void _dbg(const char *sdbg, TH h, TA... a) {
    while(*sdbg != ',') cout << *sdbg++;
    cout << '=' << (h) << ','; 
    _dbg(sdbg+1, a...);
}
#endif
\end{minted}

\subsection*{Reminders}
\added
\begin{figure}[h]
    \centering
    \includegraphics[height=0.2\textheight]{./cheer_potato.png}
    \caption{풀다가 막힐 때는 이 그림을 봅시다. 아자아자 화이팅!}
\end{figure}

\section{Math}

\subsection{Basic Mathematics}

\subsubsection{Trigonometry}
\begin{itemize}
\item $ \sin(\alpha \pm \beta) = \sin\alpha \cos\beta \pm \cos\alpha \sin\beta $
\item $ \cos(\alpha \pm \beta) = \cos\alpha \cos\beta \mp \sin\alpha \sin\beta $
\item $ \tan(\alpha \pm \beta) = \frac{\tan\alpha \pm \tan\beta}{1 \mp \tan\alpha \tan\beta}$
\item $ \sin 2\theta = 2\sin\theta\cos\theta $
\item $ c^2 = a^2 + b^2 - 2ab\cos\gamma$
\end{itemize}
\subsubsection{Generating Function}
\begin{itemize}
    \item $ \sum_{n} (pn + q)x^n = \frac{p}{1-x} + \frac{q}{(1-x)^2}$ (Arithmetic progression)
    \item $ \sum_{n} (rx)^n = (1 - rx)^{-1}$ (Geometric progression)
    \item $ \sum_{n} \binom{m}{n} x^n = (1+x)^m$ (Binomial coefficient)
    \item $ \sum_{n} \binom{m+n-1}{n} x^n = (1-x)^{-m}$ (Multiset coefficient)
\end{itemize}
\subsubsection{Calculus}
\begin{itemize}
    \item $\int_{a}^{b} f(x)\ dx = \frac{b-a}{6}\left[f(a) + 4f(\frac{a+b}{2}) + f(b)\right]$ (Simpson's Rule)
    \item $\int u'v\ dx = uv - \int uv'\ dx$ (Integration by parts)
\end{itemize}
\subsubsection{Gaussian Elimination}
\added...?
\subsection{Number Theory}
\subsubsection{Lattice Points under Line}
\begin{minted}{cpp}
// 0 <= x < n, 0 < y <= (a/c)x+(b/c)
ll calc(ll a,ll b,ll c,ll n){
    if(!n)return 0;
    ll tmp=a/c*n*(n-1)/2;
    tmp+=b/c*n;
    return tmp+calc(c,(a*n+b)%c,a%c,((a%c)*n+b%c)/c);
}
\end{minted}
\subsubsection{Shanks' Baby-step Giant-step}
\revised
\begin{minted}{cpp}
ll mexp(ll x, ll y, ll p) {
    if(!y) return 1;
    if(y & 1) return x * mexp(x*x%p, y>>1, p) % p;
    return mexp(x*x%p, y>>1, p);
}
vector<ll> get_factor(ll n) {
    vector<ll> v;
    for(ll i=2;i*i<=n;i++) {
        if(n % i == 0) {
            v.push_back(i);
            while(n % i == 0) n /= i;
        }
    }
    if(n > 1) v.push_back(n);
    return v;
}
ll get_primitive(ll n) {
    ll phi = n-1; // assume n is prime
    vector<ll> fact = get_factor(phi);
    for(ll x=2;x<=n;x++) {
        int yes = 1;
        for(ll y : fact) {
            yes &= (mexp(x, phi / y, n) != 1);
        }
        if(yes) return x;
    }
    return -1;
}
// find x s.t. x^k mod n = a -> (g^k)^y mod n = a, where x = g^y
ll bsgs(ll k, ll a, ll n) {
    ll g = get_primitive(n);
    ll phi = n-1; // assume n is prime
    if(g == -1) return -1;
    ll m = ceil(sqrt(n) + 1e-9);
    vector<pl> prec(m);
    for(ll j=0;j<m;j++) {
        prec[j] = {mexp(g, j * k % phi, n), j};
    }
    sort(prec.begin(), prec.end());
    ll cur = a, ncur = mexp(g, (phi - m) * k % phi, n);
    for(ll i=0;i<m;i++) {
        auto it = lower_bound(prec.begin(), prec.end(), pl(cur, 0));
        if(it->first == cur) {
            ll ans = mexp(g, (i*m + it->second) % phi, n);
            assert(mexp(ans, k, n) == a);
            return ans;
        }
        cur = cur * ncur % n;
    }
    return 0;
}
\end{minted}
\subsubsection{Extended Euclidean Algorithm}
\begin{minted}{cpp}
// ax + by = gcd(a,b) => x ? y ?
typedef long long int ll
pair<ll,ll> ext_gcd(ll a,ll b) {
    if(b) {
        auto tmp = ext_gcd(b, a%b);
        return {tmp.second, tmp.first - (a/b) * tmp.second};
    }
    else return {1, 0};
}

// ax = 1 mod M, x?
ll mod_inv(ll a, ll M) {
    return (ext_gcd(a, M).first + M) % M;
}
\end{minted}
\subsubsection{Chinese Remainder Theorem}
\begin{minted}{cpp}
typedef long long int ll;
ll CRT(vector<ll> rem, vector<ll> mod, int k) {
    ll m = 1;
    for(auto i : mod) m *= i;
    ll ret = 0;
    for(int i = 0 ; i < k ; i++) {
        ll tmp = (m / mod[i]) % mod[i];
        ll si = mod_inv(tmp, mod[i]);
        ret += (rem[i] * si % m) * (m / mod[i]) % m;
        ret %= m;
    }
    return ret;
}
\end{minted}
$ x \equiv a \pmod m$, $ x \equiv b \pmod n $일 경우 해가 존재하려면 $ a \equiv b \mod \gcd(m, n) $이여야 함.

$ g = \gcd(n, m) = um + vn$이라 할 때 $ x \equiv (avn + bum) / g \pmod \text{lcm}(n, m)$.
\added
\subsubsection{M\"obius Inversion Formula}
$ \forall n \in \mathbb{N}\ g(n) = \sum_{d|n} f(d) \implies f(n) = \sum_{d|n} \mu(d) g(n/d) $

\subsection{FFT}
FFT : $
(a_0, a_1, \cdots, a_{n-1}) \mapsto (\sum_{j=0}^{n-1}a_0 (\omega^0)^j, \sum_{j=0}^{n-1}a_1 (\omega^1)^j, \cdots, \sum_{j=0}^{n-1}a_{n-1} (\omega^{n-1})^j)
$
\begin{minted}{cpp}
void fft(vector<base>& a, bool inv) {
    int n = a.size(), j = 0;
    vector<ll> roots(n/2);
    for(int i=1;i<n;i++) {
        int bit = (n >> 1);
        while(j >= bit) {
            j -= bit;
            bit >>= 1;
        }
        j += bit;
        if(i < j) swap(a[i], a[j]);
    }

    double ang = 2 * acos(-1) / n * (inv ? -1 : 1);
    for(int i=0;i<n/2;i++) {
        roots[i] = base(cos(ang * i), sin(ang * i));
    }

    /* In NTT, let prr = primitive root. Then,
    int ang = mexp(prr, (mod - 1) / n);
    if(inv) ang = mexp(ang, mod - 2);
    for(int i=0; i<n/2; i++){
        roots[i] = (i ? (1ll * roots[i-1] * ang % mod) : 1);
    }
    also, make sure to apply modulus under here
    */
    for(int i=2;i<=n;i<<=1) {
        int step = n / i;
        for(int j=0;j<n;j+=i) {
            for(int k=0;k<i/2;k++) {
                ll u = a[j+k], v = a[j+k+i/2] * roots[step * k];
                a[j+k] = u+v;
                a[j+k+i/2] = u-v;
            }
        }
    }
    if(inv) for(int i=0;i<n;i++) a[i] /= n;
}

void conv(vector<base>& x, vector<base>& y) {
    int n = 2; while(n < max(x.size(), y.size())) n <<= 1;
    n <<= 1;
    x.resize(n); y.resize(n);
    fft(x, false); fft(y, false);
    for(int i=0;i<n;i++) x[i] *= y[i];
    fft(x, true); // access (ll)round(x[i].real())
}
\end{minted}

\subsection{Miller-Rabin + Pollard-Rho}
\begin{minted}{cpp}
//Prove By Solving - https://www.acmicpc.net/problem/4149
namespace miller_rabin{
    lint mul(lint x, lint y, lint mod){ return (__int128) x * y % mod; }
	lint ipow(lint x, lint y, lint p){
		lint ret = 1, piv = x % p;
		while(y){
			if(y&1) ret = mul(ret, piv, p);
			piv = mul(piv, piv, p);
			y >>= 1;
		}
		return ret;
	}
	bool miller_rabin(lint x, lint a){
		if(x % a == 0) return 0;
		lint d = x - 1;
		while(1){
			lint tmp = ipow(a, d, x);
			if(d&1) return (tmp != 1 && tmp != x-1);
			else if(tmp == x-1) return 0;
			d >>= 1;
		}
	}
	bool isprime(lint x){
		for(auto &i : {2, 3, 5, 7, 11, 13, 17, 19, 23, 29, 31, 37}){
			if(x == i) return 1;
			if(x > 40 && miller_rabin(x, i)) return 0;
		}
		if(x <= 40) return 0;
		return 1;
	}
}

namespace pollard_rho{
	lint f(lint x, lint n, lint c){
		return (c + miller_rabin::mul(x, x, n)) % n;
	}
	void rec(lint n, vector<lint> &v){
		if(n == 1) return;
		if(n % 2 == 0){
			v.push_back(2);
			rec(n/2, v);
			return;
		}
		if(miller_rabin::isprime(n)){
			v.push_back(n);
			return;
		}
		lint a, b, c;
		while(1){
			a = rand() % (n-2) + 2;
			b = a;
			c = rand() % 20 + 1;
			do{
				a = f(a, n, c);
				b = f(f(b, n, c), n, c);
			}while(gcd(abs(a-b), n) == 1);
			if(a != b) break;
		}
		lint x = gcd(abs(a-b), n);
		rec(x, v);
		rec(n/x, v);
	}
	vector<lint> factorize(lint n){
		vector<lint> ret;
		rec(n, ret);
		sort(ret.begin(), ret.end());
		return ret;
	}
};
\end{minted}
\section{Geometry}
\subsection{\texttt{struct} Point}
\begin{minted}{cpp}
const double eps = 1e-10;
template <class T>
struct point{
    typedef point P;
    T x, y;
    point(T x=0, T y=0) : x(x), y(y) {}
    bool operator< (P a) const  {return fabs(x-a.x) > eps ? x<a.x : y<a.y;}
    bool operator== (P a) const {return max(fabs(x-a.x), fabs(y-a.y)) < eps;}
    P operator+ (P a) const {return P(x+a.x, y+a.y);}
    P operator- (P a) const {return P(x-a.x, y-a.y);}
    P operator- () const {return P(-x, -y);};
    T operator* (P a) const {return x*a.x + y*a.y;} // inner prod
    T operator/ (P a) const {return x*a.y - y*a.x;} // outer prod
    T dist2() const {return x*x + y*y;}
    double dist() const {return sqrt(double(dist2()));}
    P perp() const {return P(-y, x);}; // rotate 90 deg ccw
    P mult(T t) const {return P(x*t, y*t);}
    P unit() const {return P(x/dist(), y/dist());}
    P rotate(double a){
        return P(x*cos(a)-y*sin(a), x*sin(a)+y*cos(a));
    }
};
\end{minted}
\subsubsection{Sorting Points by Angle}
\begin{minted}{cpp}
// credit : http://koosaga.com/97
auto angle_sort = [&](const point &a, const point &b){
    if((a > point(0, 0)) ^ (b > point(0, 0))) return a > b;
    if(a / b != 0) return a / b > 0;
    return a.dist() < b.dist(); // norm
}; // clockwise sort
\end{minted}
\subsection{Intersections}
벡터 내적 외적 이용하면 생각보다 간단 / 선분은 예외처리가 좀 많음
\added
\subsubsection{Line-Line intersection}

\subsubsection{Line-Segment intersection}

\subsubsection{Segment-Segment Intersection}

\subsubsection{Circle-Line Intersection}

\subsection{Projection, Reflection}
\added
\subsection{Convex Hull}
\begin{minted}{cpp}
vector<pll> get_CV(vector<pll> V){
    sort(V.begin(), V.end());
    sort(V.begin() + 1, V.end(), [&](pll x, pll y){
        pll xx = x - V[0];
        pll yy = y - V[0];
        ll res = xx / yy;
        if(res != 0) return res > 0;
        if(xx.first != yy.first) return xx.first < yy.first;
        return xx.second < yy.second;
    });

    vector<pll> ret;
    for(auto val : V){
        while(ret.size() > 1){
            pll xx = ret[ret.size() - 2] - val;
            pll yy = ret[ret.size() - 1] - val;

            if(xx / yy <= 0) ret.pop_back();
            else break;
        }
        ret.push_back(val);
    }

    return ret;
}
\end{minted}
\subsection{Rotating Calipers}
\begin{minted}{cpp}
void rotating_calipers(vector<pll> CV){
    int pos = 0;
    for(int i = 0 ; i < CV.size() ; i++) if(CV[pos] < CV[i]) pos = i;

    int ind1 = 0, ind2 = pos;
    ll dist = (CV[ind1] - CV[ind2]) * (CV[ind1] - CV[ind2]);

    auto get_v = [&](int x) { return CV[(x + 1) % CV.size()] - CV[x];};
    for(int i = 0 ; i < CV.size() ; i++){
        pll v = get_v(i);
        while((-v) / get_v(pos) < 0) pos = (pos + 1) % CV.size();
        ll tmp_dist = (CV[pos] - CV[i]) * (CV[pos] - CV[i]);
        if(dist < tmp_dist) {
            dist = tmp_dist;
            ind1 = i;   ind2 = pos;
        }
    }

    printf("%lld %lld %lld %lld\n", CV[ind1].first, CV[ind1].second, CV[ind2].first, CV[ind2].second);
}
\end{minted}

\subsection{Smallest Enclosing Circle}
\begin{minted}{cpp}
//Prove By Solving - https://www.acmicpc.net/problem/11930
int main(){
    scanf("%d", &N);

    for(int i = 1 ; i <= N ; i++) scanf("%lf%lf%lf", &A[i].x, &A[i].y, &A[i].z);

    int t = 70000;
    double rate = 1.0;
    point cur = (point){0, 0, 0};
    for(int i = 1 ; i <= t; i++){
        int ind = 1;
        for(int j = 1 ; j <= N ; j++) if((A[j] - cur) * (A[j] - cur) > (A[ind] - cur) * (A[ind] - cur)) ind = j;
        cur = cur + (A[ind] - cur) * rate;
        rate *= 0.99;
    }

    double r = 0;
    for(int i = 1 ; i <= N ; i++) r = max(r, (A[i] - cur) * (A[i] - cur));
    cout << sqrt(r);
    return 0;
}
\end{minted}
\subsection{Polygon Area}
\added

\subsubsection{Polygon Area}

\subsubsection{Polygon Overlapping}

\section{Strings}

\subsection{Aho-Corasick Algorithm}
\begin{minted}{cpp}
namespace aho_corasick {
    const int MAXN = 100000, MAXC = 26;
    int trans[MAXN+1][MAXC];
    int fail[MAXN+1];
    bool term[MAXN+1];

    void build(const vector<string> &v) {
        setz(trans), setz(fail), setz(term);
        int cnode = 1;

        repi(s, v) {
            int p = 0;
            repi(j, s) {
                char c = j-'a';
                if (!trans[p][c]) trans[p][c] = cnode++;
                p = trans[p][c];
            }
            term[p] = true;
        }

        queue<int> q; rep(i, MAXC) if (trans[0][i]) q.push(trans[0][i]);
        while(!q.empty()) {
            int t = q.front(); q.pop();
            rep(i, MAXC) {
                if (trans[t][i]) {
                    int p = fail[t];
                    while(p and not trans[p][i]) p = fail[p];
                    p = trans[p][i];
                    fail[trans[t][i]] = p;
                    if (term[p]) term[trans[t][i]] = true;
                    q.push(trans[t][i]);
                }
            }
        }
    }

    bool query(string &t) {
        int p = 0;
        repi(i, t) {
            char c = i-'a';
            while(p and not trans[p][c]) p = fail[p];
            p = trans[p][c];
            if (term[p]) return true;

        }
        return false;
    }
}
\end{minted}

\subsection{Suffix Array}
\begin{minted}{cpp}
// str  : abracadabra
// SA   : 10 7 0 3 5 8 1 4 6 9 2
// LCP  : 1 4 1 1 0 3 0 0 0 2
vector<int> make_sa(const string& s) {
    int n = s.length();
    int lim = max(128, n+1);
    vector<int> sa(n), g(n+1), ng(n+1), cnt(lim), ind(lim+1);
    for(int i=0;i<n;i++) {
        sa[i] = i; g[i] = s[i];
    }
    g[n] = 0;
    for(int t=1;t<s.length();t<<=1)
    {
        auto cmp = [&] (int a, int b) {
            return g[a] != g[b] ? g[a] < g[b] : g[a+t] < g[b+t];
        };
        for(int i=0;i<n;i++)    cnt[g[min(i+t, n)]]++;
        for(int i=1;i<lim;i++)  cnt[i] += cnt[i-1];
        for(int i=n-1;i>=0;i--) ind[--cnt[g[min(i+t, n)]]] = i;
        for(int i=0;i<lim;i++)  cnt[i] = 0;
        for(int i=0;i<n;i++)    cnt[g[i]]++; // same as cnt[g[ind[i]]]++
        for(int i=1;i<lim;i++)  cnt[i] += cnt[i-1];
        for(int i=n-1;i>=0;i--) sa[--cnt[g[ind[i]]]] = ind[i];
        ng[sa[0]] = 1;
        for(int i=1;i<n;i++) {
            ng[sa[i]] = ng[sa[i-1]] + cmp(sa[i-1], sa[i]);
        }
        g = ng;
        
        fill(cnt.begin(), cnt.end(), 0);
        fill(ind.begin(), ind.end(), 0);
    }
    return sa;
}

vector<int> make_lcp(const string& s, const vector<int>& sa) {
    int n = s.length();
    vector<int> lcp(n-1), rank(n);
    for(int i=0;i<n;i++)
        rank[sa[i]] = i;
    int len = 0;
    for(int i=0;i<n;i++) {
        if(rank[i]) {
            int j = sa[rank[i]-1];
            int lc = n - max(i,j);
            while(len < lc && s[i+len] == s[j+len]) len++;
            lcp[rank[i]-1] = len;
        }
        if(len) len--;
    }
    return lcp;
}
\end{minted}
\subsection{Manacher's Algorithm}
\begin{minted}{cpp}
// 0-based
// s   = # h # e # l # l # o #
// ret = 0 1 0 1 0 1 2 1 0 1 0

vector<int> manacher(const string& s) {
    int n = s.size(), r = -1, k = -1;
    vector<int> p(n);
    for (int i=0; i<n; i++) {
        if (i<=r) p[i] = min(r-i, p[2*k-i]);
        while (i-p[i]-1>=0 and i+p[i]+1<n and s[i-p[i]-1] == s[i+p[i]+1]) p[i]++;
        if (r < i+p[i]) r = i+p[i], k = i;
    }
    return p;
}
\end{minted}
\subsection{Manacher's Algorithm}
\begin{minted}{cpp}
// 0-based
// s   = # h # e # l # l # o #
// ret = 0 1 0 1 0 1 2 1 0 1 0

vector<int> manacher(const string& s) {
    int n = s.size(), r = -1, k = -1;
    vector<int> p(n);
    for (int i=0; i<n; i++) {
        if (i<=r) p[i] = min(r-i, p[2*k-i]);
        while (i-p[i]-1>=0 and i+p[i]+1<n and s[i-p[i]-1] == s[i+p[i]+1]) p[i]++;
        if (r < i+p[i]) r = i+p[i], k = i;
    }
    return p;
}
\end{minted}

\subsection{Z Algorithm}
\begin{minted}{cpp}
// 0-based
// s   = a b c a b a b c a
// ret = 9 0 0 2 0 4 0 0 1 

vector<int> z_algo(const string &s) {
    int l = 0, r = 0, N = sz(s);
    vector<int> Z(N);
    Z[0] = N;
    repp(i, 1, N) {
        if (i > r) {
            l = r = i;
            while(r < N and s[r] == s[r-l]) r++;
            r--;
            Z[i] = r-l+1;
        } else {
            int k = i-l;
            if (Z[k] < r-i+1) Z[i] = Z[k];
            else {
                l = i;
                while(r < N and s[r] == s[r-l]) r++;
                r--;
                Z[i] = r-l+1;
            }
        }
    }
    return Z;
}
\end{minted}

\subsection{Lexicographically Smallest String Rotation}
\begin{minted}{cpp}
// rotate(v.begin(), v.begin()+min_rotation(v), v.end());
int min_rotation(string s) {
	int a=0, N=sz(s); s += s;
	rep(b,0,N) rep(i,0,N) {
		if (a+i == b || s[a+i] < s[b+i]) {b += max(0, i-1); break;}
		if (s[a+i] > s[b+i]) { a = b; break; }
	}
	return a;
}
\end{minted}

\section{Graph Theory}

\subsection{Strongly Connected Component}
\begin{minted}{cpp}
vector<int> adj[MAX_V];
int finished[MAX];
int dfsn[MAX_V];
int cnt = 1;
vector<vector<int>> SCC;
stack<int> s;

int dfs(int cur) {
    dfsn[cur] = cnt++;
    s.push(cur);
    int res = dfsn[cur];
    for (int n : adj[cur]) {
        if (!dfsn[n]) res = min(res, dfs(n));
        else if (!finished[n]) res = min(res, dfsn[n]);
    }
    if (res == dfsn[cur]) {
        vector<int> sub;
        int t = -1;
        do {
            t = s.top(); s.pop();
            finished[t] = 1;
            sub.push_back(t);
        } while (t != cur);
        SCC.push_back(sub);
    }
    return res;
}
\end{minted}
\subsection{2-SAT}
\added
\subsection{Biconnected Component}
\added
\subsection{Euler Tour}
\begin{minted}{cpp}
struct Edge{
    int to, cnt; // to: 인접한 정점, cnt: 남은 사용 횟수
    Edge *dual; // dual: 역방향 간선을 가리키는 포인터
    Edge(): Edge(-1, 0){}
    Edge(int to1, int cnt1): to(to1), cnt(cnt1), dual(nullptr){}
};
void Eulerian(int curr){
    for(Edge *e: adj[curr]){
        if(e->cnt > 0){
            e->cnt--;
            e->dual->cnt--;
            Eulerian(e->to); // dfs
        }
    }
    cout << curr << '\n';
}
\end{minted}
\subsection{Offline Dynamic Connectivity}
\added
\subsection{Heavy-Light Decomposition}
\added % 각 서브트리에서 크기 순으로 정렬하기
\subsection{Dominator Tree}
\added
\subsection{Negative Cycle Detection}
\added
\subsection{Tree Compress}
\added

\section{Network Flow}

\subsection{Theorems}
\noindent\textbf{Max-flow Min-cut theorem} : 정점 $ s $에서 정점 $ t $까지 흐를 수 있는 최대 유량(max-flow)은 정점 $ s $와 정점 $ t $를 분리하는 간선들의 가중치 합(min-cut)과 같다.\\
\textbf{Vertex cover} : 어떤 그래프의 정점의 집합 $ S $에 대해 그래프의 모든 간선이 $ S $의 원소 중 최소 하나와 연결되어 있을 때, $ S $를 해당 그래프의 vertex cover라고 하며, minimum vertex cover는 최소 개수의 정점을 사용한 vertex cover이다.\\
\textbf{Independent set} : 어떤 그래프의 정점의 집합 $ S $에 대해 $ S $의 서로 다른 두 정점을 연결하는 간선이 없을 때, $ S $를 해당 그래프의 independent set이라고 하며, maximum independent set은 최대 개수의 정점을 사용한 independent set이다.\\
\textbf{Matching (independent edge set)} : 어떤 그래프의 간선의 집합 $ S $에 대해 $ S $의 서로 다른 두 간선이 공통된 정점을 가지지 않을 때, $ S $를 해당 그래프의 matching이라고 하며, maximum matching은 최대 개수의 간선을 사용한 matching이다.\\
\textbf{K\"onig's theorem} : 이분 그래프의 maximum matching의 크기는 minimum vertex cover의 것과 같다.\\
\textbf{Dinic's Algorithm} : 시간 복잡도 $O(V^2E)$, unit capacity에서는 $\min(V^{2/3}E, E^{3/2})$.\\
\textbf{Circulation Problem} : 새로운 source/sink $s_{n}$, $t_{n}$를 만들어서 다음과 같이 간선을 추가하고 $maxflow(s_{n} \to t_{n}) = \sum l_i$인지 확인, 이후 $s \to t$로 maxflow
\begin{itemize}
\item $s_{n} \to b\ (l)$, $a \to t_{n}\ (l)$, $a \to b\ (r-l)$, $t \to s\ (\infty)$
\end{itemize}

\subsection {Dinic's Algorithm}
\begin{minted}{cpp}
const int INF = 1e9;
struct Dinic{
    int N;
    struct edge{
        int index, cap, rev;
        edge() : index(0), cap(0), rev(0) {}
        edge(int index, int cap, int rev) : index(index), cap(cap), rev(rev) {}
    };

    vector<vector<edge>> ADJ;
    vector<int> R, W;

    Dinic() {}
    Dinic(int N) : N(N){
        ADJ.resize(N);  R.resize(N);    W.resize(N);
    }

    void CE(int node1, int node2, int cap){
        ADJ[node1].push_back(edge(node2, cap, ADJ[node2].size()));
        ADJ[node2].push_back(edge(node1, 0, ADJ[node1].size() - 1));
    }

    bool bfs(int src, int sink){
        for(int i = 0 ; i < R.size() ; i++) R[i] = -1;
        R[src] = 0;
        queue<int> Q; Q.push(src);
        while(Q.size()){
            int here = Q.front();   Q.pop();
            for(auto e : ADJ[here]){
                if(e.cap > 0 && R[e.index] == -1) R[e.index] = R[here] + 1, Q.push(e.index);
            }
        }
        return R[sink] != -1;
    }

    int dfs(int here, int sink, int f){
        if(here == sink) return f;
        for(int &i = W[here] ; i < ADJ[here].size() ; i++){
            auto &e = ADJ[here][i];
            if(e.cap > 0 && R[here] < R[e.index]){
                int res = dfs(e.index, sink, min(f, e.cap));
                if(res) {
                    e.cap -= res;
                    ADJ[e.index][e.rev].cap += res;
                    return res; 
                }
            }
        }
        return 0;
    }

    int solve(int src, int sink){
        int ret = 0;
        while(bfs(src, sink)){
            for(int i = 0 ; i < N ; i++) W[i] = 0;
            int res;
            while((res = dfs(src, sink, INF))) ret += res;
        }
        return ret;
    }
};
\end{minted}
\subsection{MCMF with SPFA}
\begin{minted}{cpp}
const int INF = 1e9;
struct MCMF {
	struct EDGE {
		int there, cap, cost, rev;

		EDGE() : there(0), cap(0), cost(0), rev(0) {}
		EDGE(int there, int cap, int cost, int rev) : there(there), cap(cap), cost(cost), rev(rev) {}
	};

	int N;
	vector<vector<EDGE>> ADJ;
	vector<int> R, INQ, C, I;

	MCMF() : N(0) {}
	MCMF(int N) : N(N) { ADJ.resize(N + 1); R.resize(N + 1); INQ.resize(N + 1);	 C.resize(N + 1); I.resize(N + 1); }

	void connect_edge(int i, int j, int cap, int cost) {
		ADJ[i].push_back(EDGE(j, cap, cost, ADJ[j].size()));
		ADJ[j].push_back(EDGE(i, 0, -cost, ADJ[i].size() - 1));
	}
	bool SPFA(int src, int sink) {
		queue<int> Q;	Q.push(src);
		fill(R.begin(), R.end(), -1);	R[src] = 0;
		fill(C.begin(), C.end(), -1);	C[src] = 0;
		fill(INQ.begin(), INQ.end(), 0);	INQ[src] = 1;
		while (Q.size()) {
			int here = Q.front();	Q.pop();
			INQ[here] = 0;
			for (int i = 0; i < ADJ[here].size(); i++) {
				auto e = ADJ[here][i];
				if (e.cap > 0 && (C[e.there] == -1 || C[e.there] > C[here] + e.cost)) {
					C[e.there] = C[here] + e.cost;	R[e.there] = here;	I[e.there] = i;
					if (!INQ[e.there]) INQ[e.there] = 1, Q.push(e.there);
				}
			}
		}
		if (C[sink] == -1) return false;
		return true;
	}

	pii mcmf(int src, int sink) {
		pii ret = { 0, 0 };
		while (SPFA(src, sink)) {
			int flow = INF, cost = 0;
			for (int here = sink; here != src; here = R[here]) flow = min(flow, ADJ[R[here]][I[here]].cap);
			for (int here = sink; here != src; here = R[here]) {
				auto &e = ADJ[R[here]][I[here]];
				cost += e.cost * flow;
				e.cap -= flow;
				ADJ[e.there][e.rev].cap += flow;
			}
			ret.first += flow, ret.second += cost;
		}
		return ret;
	}
};
\end{minted}

\subsection{Hungarian Method}
\begin{minted}{cpp}
namespace Hung {
    const int MX = 2000;
    // IMPORTANT : n <= m, 1-based
    using T = long double;

    T maxv = 1e200;
    T a[MX][MX], n, m;

    void init(int nn, int mm) { n = nn; m = mm; }
    void set_value(int x, int y, T val) { a[x][y] = val; }
    T solve(vector <int> &ans) {
        vector<T> v(m+1), u(n+1);
        vector<int> p (m+1), way (m+1);
        for (int i=1; i<=n; ++i) {
            p[0] = i;
            int j0 = 0;
            vector<T> minv (m+1, maxv);
            vector<char> used (m+1, false);
            do {
                used[j0] = true;
                T delta = maxv;
                int i0 = p[j0], j1;
                for (int j=1; j<=m; ++j) if (!used[j]) {
                    T cur = a[i0][j]-u[i0]-v[j];
                    if (cur < minv[j]) {
                        minv[j] = cur, way[j] = j0;
                    }
                    if (minv[j] < delta) {
                        delta = minv[j], j1 = j;
                    }
                }
                for (int j=0; j<=m; ++j) {
                    if (used[j]) {
                        u[p[j]] += delta, v[j] -= delta;
                    }
                    else {
                        minv[j] -= delta;
                    }
                }
                j0 = j1;
            } while (p[j0] != 0);
            do {
                int j1 = way[j0];
                p[j0] = p[j1];
                j0 = j1;
            } while (j0);
        }
        ans.resize(n + 1);
        for(int j=1;j<=m;++j) {
            ans[p[j]] = j;
        }
        return -v[0];
    }
}
\end{minted}
\subsection{Hopcroft-Karp Algorithm}
\begin{minted}{cpp}
struct hopcroft_karp{
    int N;
    vector<vector<int>> ADJ;
    vector<int> L, rev, used;

    hopcroft_karp() {}
    hopcroft_karp(int N) : N(N) {
        ADJ.resize(N);  
        L.resize(N), rev.resize(N, -1), used.resize(N, 0);
    }

    void CE(int here, int there){
        ADJ[here].push_back(there);
    }

    void bfs(){
        queue<int> Q;
        for(int i = 0 ; i < N ; i++) {
            if(used[i]) L[i] = -1;
            else L[i] = 0, Q.push(i);
        }

        while(Q.size()){
            int here = Q.front();   Q.pop();
            for(int there : ADJ[here]){
                if(rev[there] != -1 && L[rev[there]] == -1) {
                    L[rev[there]] = L[here] + 1;
                    Q.push(rev[there]);
                }
            }
        }
    }

    bool dfs(int here){
        for(int there : ADJ[here]){
            if(rev[there] == -1 || (L[here] < L[rev[there]] && dfs(rev[there]))){
                rev[there] = here;
                used[here] = 1;
                return true;
            }
        }
        return false;
    }

    int solve(){
        int ret = 0;
        while(1){
            bfs();
            int res = 0;
            for(int i = 0 ; i < N ; i++) {
                if(used[i]) continue;
                res += dfs(i);
            }
            if(res == 0) break;
            ret += res;
        }
        return ret;
    }
};
\end{minted}

\section{Optimization Tricks}

\subsection{Knuth Optimization}
\begin{itemize}
\item Recurrence : $ D[i][j] = \min_{i<k<j}(D[i][k] + D[k][j]) + C[i][j]$
\item Quadrangle Inequality : $ C[a][c] + C[b][d] \leq C[a][d] + C[b][c],\ a \leq b \leq c \leq d$
\item Monotonicity : $ C[b][c] \leq C[a][d],\ a \leq b \leq c \leq d$
\item $ A[i][j] = $ (min. $ k $ \textit{s.t.} $ D[i][j] $ is min.). Then $ A[i][j-1] \leq A[i][j] \leq A[i+1][j]$
\item $ O(N^2) $ time complexity
\end{itemize}
\begin{minted}{cpp}
// opt[i-1][i] = i
for(int d=2;d<=n;d++) {
    for(int i=1;i+d<=n+1;i++) {
        for(int k=opt[i][j-1], j=i+d; k<=opt[i+1][j]; k++) {
            int v = dp[i][k] + dp[k][j] + c[i][j];
            if(dp[i][j] > v) dp[i][j] = v, opt[i][j] = k;
        }
    }
}
\end{minted}

\subsection{Divide and Conquer Optimization}
\begin{itemize}
\item Recurrence : $ D[t][i] = \min_{k<i} (D[t-1][k] + C[k][i]) $
\item Min index : $ A[t][i] \leq A[t][i+1] $ ($ A[t][i] = $ (min. $ k $ \textit{s.t.} $ D[t][i] $ is min.))
\subitem[-] Quadrangle Inequality : $ C[a][c] + C[b][d] \leq C[a][d] + C[b][c],\ a \leq b \leq c \leq d$
\item Able to Divide and Conquer base on calculating $ D[t][i] $
\item $ O(T N \lg N) $ time complexity
\end{itemize}
\begin{minted}{cpp}
// range of index : [l,r]
// range of dp    : [s,e]
void dnc(int t, int l, int r, int s, int e)
{
    if(s > e) return;
    int m = (s+e)/2;
    D[t][m] = 2e9;
    for(int k=l;k<m&&k<=r;k++)
    {
        int tmp = D[t-1][k] + C[k][m];
        if(D[t][m] > tmp)
            D[t][m] = tmp, A[t][m] = k;
    }
    dnc(t, l, A[t][m], s, D[t][m]-1);
    dnc(t, A[t][m], r, D[t][m]+1, e);
}
\end{minted}
\subsection{Convex Hull Trick}
\begin{itemize}
\item Recurrence : $ dp[i] = \min_{j<i} (dp[j] + a[i]b[j]),\ b[i-1] \leq b[i] $
\item Think as $ dp[x = a[i]] = \min_{j<i}(b[j] \cdot x + dp[j]) $
\item Thus push lines and find minimum (by binary search)
\item If $ a[i] \leq a[i+1] $ sweeping is possible
\item Intersection of $ y = a_ix +b_i $ and $ y = a_{i+1}x + b_{i+1} $ : $ x = \frac{b_{i+1}-b_i}{a_i-a_{i+1}} $
\end{itemize}

\subsection{Centroid Decomposition}
\begin{minted}{cpp}
// credit : https://gist.github.com/igorcarpanese/75162f3253bd230abd0f32f9950bf384
int dfs(int u, int p) {
    sub[u] = 1;
    for (auto v : tree[u])
        if (v != p) sub[u] += dfs(v, u);
    return sub[u] + 1;
}

int centroid(int u, int p, int r) { // r : root
    for (auto v : tree[u])
        if (v != p and sub[v] > sub[r]/2) return centroid(v, u);
    return u;
}
\end{minted}

%\subsection{IOI 2016 Alien Trick}
%\added...? (move to misc?)

\section{Data Structure}

\subsection{Persistent Segment Tree}
\begin{minted}{cpp}
const MAXN = 1e5 + 10;
struct node{
    node *l, *r;
    int cnt;
    node () {}  
} pool[(1 << 17) * 17], *tree_head[MAXN];

int tcnt;
node* alloc(){
    memset(pool + tcnt, 0, sizeof(node));
    return pool + tcnt++;
}

node * init(int l, int r){
    node *ret = alloc();
    if(l != r) {
        int mid = (l + r) / 2;
        ret->l = init(l, mid);
        ret->r = init(mid + 1, r);
    }
    return ret;
}

void update(node * here, node *par, int l, int r, int val){
    if(l == r) {
        here->cnt = par->cnt + 1;
        return;
    }

    int mid = (l + r) / 2;
    if(val <= mid){
        here->l = alloc();
        here->r = par->r;
        update(here->l, par->l, l, mid, val);
    }
    else {
        here->l = par->l;
        here->r = alloc();
        update(here->r, par->r, mid + 1, r, val);
    }
    here->cnt = here->l->cnt + here->r->cnt;
}

int query(node *node1, node *node2, int l, int r, int k){
    if(l == r) return l;
    int ccc = node1->l->cnt - node2->l->cnt;
    int mid = (l + r) / 2;
    if(k <= ccc) return query(node1->l, node2->l, l, mid, k);
    else return query(node1->r, node2->r, mid + 1, r, k - ccc);
}
\end{minted}

\subsection{Link-Cut Tree}
\begin{minted}{cpp}
struct node{
    node *pp, *p, *l, *r;
    int val;
    node(){ p = 0, l = 0, r = 0;}
    node(int val) : val(val) { p = 0, l = 0, r = 0;}
};

void push(node *x){}
void pull(node *x){}

void rotate(node *x){
	if(!x->p) return;
	push(x->p);  // if there's lazy stuff
	push(x);
	node *p = x->p;
	bool is_left = (p->l == x);
	node *b = (is_left ? x->r : x->l);
	x->p = p->p;
	if(x->p && x->p->l == p) x->p->l = x;
	if(x->p && x->p->r == p) x->p->r = x;
	if(is_left){
		if(b) b->p = p;
		p->l = b;
		p->p = x;
		x->r = p;
	}
	else{
		if(b) b->p = p;
		p->r = b;
		p->p = x;
		x->l = p;
	}
	pull(p); // if there's something to pull up
	pull(x);
	//if(!x->p) root = x; // IF YOU ARE SPLAY TREE
	if(p->pp){ // IF YOU ARE LINK CUT TREE
		x->pp = p->pp;
		p->pp = nullptr;
	}
}
void splay(node *x){
	while(x->p){
		node *p = x->p;
		node *g = p->p;
		if(g){
			if((p->l == x) ^ (g->l == p)) rotate(x);
			else rotate(p);
		}
		rotate(x);
	}
}
void access(node *x){
	splay(x);
	push(x);
	if(x->r){
		x->r->pp = x;
		x->r->p = nullptr;
		x->r = nullptr;
	}
	pull(x);
	while(x->pp){
		node *nxt = x->pp;
		splay(nxt);
		push(nxt);
		if(nxt->r){
			nxt->r->pp = nxt;
			nxt->r->p = nullptr;
			nxt->r = nullptr;
		}
		nxt->r = x;
		x->p = nxt;
		x->pp = nullptr;
		pull(nxt);
		splay(x);
	}
}
node *root(node *x){
	access(x);
	while(x->l){
		push(x);
		x = x->l;
	}
	access(x);
	return x;
}
node *par(node *x){
	access(x);
	if(!x->l) return nullptr;
	push(x);
	x = x->l;
	while(x->r){
		push(x);
		x = x->r;
	}
	access(x);
	return x;
}
node *lca(node *s, node *t){
	access(s);
	access(t);
	splay(s);
	if(s->pp == nullptr) return s;
	return s->pp;
}
void link(node *par, node *son){
	access(par);
	access(son);
	//son->rev ^= 1;  // remove if needed
	push(son);
	son->l = par;
	par->p = son;
	pull(son);
}
void cut(node *p){
	access(p);
	push(p);
	if(p->l){
		p->l->p = nullptr;
		p->l = nullptr;
	}
	pull(p);
}
\end{minted}
\subsection{Li-Chao Tree}
\added...?
\subsection{Dynamic Convex Hull}
\begin{minted}{cpp}
// https://github.com/niklasb/contest-algos/blob/master/convex_hull/dynamic.cpp
const ll is_query = -(1LL<<62);
struct Line {
    ll m, b;
    mutable function<const Line*()> succ;
    bool operator<(const Line& rhs) const {
        if (rhs.b != is_query) return m < rhs.m;
        const Line* s = succ();
        if (!s) return 0;
        ll x = rhs.m;
        return b - s->b < (s->m - m) * x;
    }
};
struct HullDynamic : public multiset<Line> { // will maintain upper hull for maximum
    bool bad(iterator y) {
        auto z = next(y);
        if (y == begin()) {
            if (z == end()) return 0;
            return y->m == z->m && y->b <= z->b;
        }
        auto x = prev(y);
        if (z == end()) return y->m == x->m && y->b <= x->b;
        return (x->b - y->b)*(z->m - y->m) >= (y->b - z->b)*(y->m - x->m);
    }
    void insert_line(ll m, ll b) {
        auto y = insert({ m, b });
        y->succ = [=] { return next(y) == end() ? 0 : &*next(y); };
        if (bad(y)) { erase(y); return; }
        while (next(y) != end() && bad(next(y))) erase(next(y));
        while (y != begin() && bad(prev(y))) erase(prev(y));
    }
    ll eval(ll x) {
        auto l = *lower_bound((Line) { x, is_query });
        return l.m * x + l.b;
    }
};
\end{minted}
\subsection{Stern-Brocot Tree}
\begin{minted}{cpp}
// __int128 is recommended
bool test(ll a, ll b) { // for testing directions, vary by prob
    // return true if (true value) >= a/b
    ll n = 0, m = 1;

    rep(i, N) {
        if (n < m*A[i].fi) n = A[i].fi, m = 1;

        ll c = b*n+m*a, d = m*b;
        ll g = gcd(c, d);
        n = c/g;
        m = d/g;

        if (n > m*A[i].se) return false;
    }
    return true;
}

pair<ll, ll> stern_brocot(ll M, ll N) {
    // M : max value
    // N : max divisor
    // if result is a/b, return as {a, b}
    ll a = 0, b = 1; // l
    ll c = 1, d = 0; // r
    int l, r;
    bool chg = true;

    while(chg) {
        chg = false;

        // to left
        l = 0, r = (N-d-1)/b+1;
        while(l < r) {
            int mid = (l+r+1)/2;
            if (test(a*mid+c, b*mid+d)) r = mid-1;
            else l = mid;
        }

        c += a*l;
        d += b*l;
        chg |= (l > 0);

        // to right
        l = 0, r = (d?(N-b-1)/d+1:M);
        while(l < r) {
            int mid = (l+r+1)/2;
            if (test(a+mid*c, b+mid*d)) l = mid;
            else r = mid-1;
        }

        a += c*l;
        b += d*l;
        chg |= (l > 0);
    }

    return {a, b};
}
\end{minted}
%\subsection{Treap}
%\added...?
\subsection{Rope}
\added
\subsection{Policy Based Data Structure}
\begin{minted}{cpp}
#include <bits/stdc++.h>
#include <ext/pb_ds/assoc_container.hpp> 
#include <ext/pb_ds/tree_policy.hpp>
#include <ext/pb_ds/detail/standard_policies.hpp>

using namespace std;
using namespace __gnu_pbds;

typedef tree<
int, 
null_type, 
less<int>, 
rb_tree_tag,
tree_order_statistics_node_update >
ordered_set;

// less<int> : not allow for duplicate
// less_equal<int> : allow for duplicate
// use upper_bound when you erase from set used less_equal

int N;

int main(void) {
    iostream::sync_with_stdio(false);
    cin.tie(nullptr);
    
    ordered_set X;

    X.insert(1);
    X.insert(2);
    X.insert(4);
    X.insert(8);
    X.insert(16);

    cout<<*X.find_by_order(1)<<endl; // 2
    cout<<*X.find_by_order(2)<<endl; // 4
    cout<<*X.find_by_order(4)<<endl; // 16
    cout<<(end(X)==X.find_by_order(6))<<endl; // true

    cout<<X.order_of_key(-5)<<endl;  // 0
    cout<<X.order_of_key(1)<<endl;   // 0
    cout<<X.order_of_key(3)<<endl;   // 2
    cout<<X.order_of_key(4)<<endl;   // 2
    cout<<X.order_of_key(400)<<endl; // 5   
}
\end{minted}
\section{Miscellaneous}
\subsection{Misc Formulae and Algorithms}
\subsubsection{Faulhaber's Formula}
$$
T(n, k) = \sum_{i=1}^{n} i^k = \dfrac{(n+1)^{k+1} - 1^{k+1} - \sum_{j=0}^{k-1}\binom{k+1}{j}T(n, j)}{\binom{k+1}{k}}
$$
Also use 
$$
(x+1) ^ d - x ^ d = 1 + \dbinom{d}{1} x + \dbinom{d}{2} x^2 + \cdots + \dbinom{d}{d-1} x^{d-1}
$$
to get each coef.
\subsubsection{Maximum Clique}
\begin{minted}{cpp}
typedef long long ll;
ll G[40]; // 0-index
int N, M;
int cur;
void get_clique(int R = 0, ll P = (1ll<<N)-1, ll X = 0){
    if((P|X) == 0){
        cur = max(cur, R);
        return;
    }
    int u = __builtin_ctzll(P|X);
    ll c = P&~G[u];
    while(c){
        int v = __builtin_ctzll(c);
        get_clique(R + 1, P&G[v], X&G[v]);
        P ^= 1ll << v;
        X |= 1ll << v;
        c ^= 1ll << v;
    }
}
\end{minted}
\subsubsection{De Brujin Sequence}
\added...?

\subsection {Highly Composite Numbers, Large Prime}
\begin{minted}{cpp}
	< 10^k          number     divisors   2 3 5 71113171923293137
	-------------------------------------------------------------
	1                    6            4   1 1
	2                   60           12   2 1 1
	3                  840           32   3 1 1 1
	4                 7560           64   3 3 1 1
	5                83160          128   3 3 1 1 1
	6               720720          240   4 2 1 1 1 1
	7              8648640          448   6 3 1 1 1 1
	8             73513440          768   5 3 1 1 1 1 1
	9            735134400         1344   6 3 2 1 1 1 1
	10          6983776800         2304   5 3 2 1 1 1 1 1
	11         97772875200         4032   6 3 2 2 1 1 1 1
	12        963761198400         6720   6 4 2 1 1 1 1 1 1
	13       9316358251200        10752   6 3 2 1 1 1 1 1 1 1
	14      97821761637600        17280   5 4 2 2 1 1 1 1 1 1
	15     866421317361600        26880   6 4 2 1 1 1 1 1 1 1 1
	16    8086598962041600        41472   8 3 2 2 1 1 1 1 1 1 1
	17   74801040398884800        64512   6 3 2 2 1 1 1 1 1 1 1 1
	18  897612484786617600       103680   8 4 2 2 1 1 1 1 1 1 1 1
	
	< 10^k    prime   # of prime          < 10^k            prime
	-------------------------------------------------------------
	1             7            4          10           9999999967  
	2            97           25          11          99999999977
	3           997          168          12         999999999989
	4          9973         1229          13        9999999999971
	5         99991         9592          14       99999999999973
	6        999983        78498          15      999999999999989
	7       9999991       664579          16     9999999999999937
	8      99999989      5761455          17    99999999999999997
	9     999999937     50847534          18   999999999999999989
\end{minted}
NTT Prime:

$469762049 = 7 \times 2^{26} + 1$. Primitive root : $3$.

$998244353 = 119 \times 2^{23} + 1$. Primitive root: $3$.

$985661441 = 235 \times 2^{22} + 1$. Primitive root: $3$.

$1012924417 = 483 \times 2^{21} + 1$. Primitive root: $5$.

\subsection {Fast Integer IO}
\begin{minted}{cpp}
// credit : https://github.com/koosaga/DeobureoMinkyuParty/blob/master/teamnote.tex
static char buf[1 << 19]; // size : any number geq than 1024 
static int idx = 0;
static int bytes = 0;
static inline int _read() {
	if (!bytes || idx == bytes) {
		bytes = (int)fread(buf, sizeof(buf[0]), sizeof(buf), stdin);
		idx = 0;
	}
	return buf[idx++];
}
static inline int _readInt() {
	int x = 0, s = 1;
	int c = _read();
	while (c <= 32) c = _read();
	if (c == '-') s = -1, c = _read();
	while (c > 32) x = 10 * x + (c - '0'), c = _read();
	if (s < 0) x = -x;
	return x;
}
\end{minted}

\subsection {C++ Tips / Environments}
\revised (with random, \verb|chrono|)
\begin{minted}{cpp}
#include <bits/stdc++.h> // magic header
using namespace std;     // magic namespace

struct StupidGCCCantEvenCompileThisSimpleCode{
	pair<int, int> array[1000000];
}; // https://gcc.gnu.org/bugzilla/show_bug.cgi?id=68203

// how to use rand (in 2017)
mt19937 rng(0xdeadbeef);
mt19937 rng(chrono::steady_clock::now().time_since_epoch().count());
int randint(int lb, int ub){ return uniform_int_distribution<int>(lb, ub)(rng); }
shuffle(permutation.begin(), permutation.end(), rng);
mt19937_64 _R(chrono::steady_clock::now().time_since_epoch().count()); // _R()

// comparator overload
auto cmp = [](seg a, seg b){return a.func() < b.func(); };
set<seg, decltype(cmp)> s(cmp); 
map<seg, int, decltype(cmp)> mp(cmp);
priority_queue<seg, vector<seg>, decltype(cmp)> pq(cmp); // max heap

// hash func overload
struct point{
int x, y;
bool operator==(const point &p)const{ return x == p.x && y == p.y; }
};
struct hasher {
size_t operator()(const point &p)const{ return p.x * 2 + p.y * 3; }
};
unordered_map<point, int, hasher> hsh;

// c++ setprecision example
#include <iostream>     // std::cout, std::fixed
#include <iomanip>      // std::setprecision

int main () {
  double f =3.14159;
  std::cout << std::setprecision(5) << f << '\n'; // 3.1416
  std::cout << std::setprecision(9) << f << '\n'; // 3.14159
  std::cout << std::fixed;
  std::cout << std::setprecision(5) << f << '\n'; // 3.14159
  std::cout << std::setprecision(9) << f << '\n'; // 3.141590000
  return 0;
}
\end{minted}

%\subsection{Some Mathematics?}
%Xudyu sieve? Cayley theorem? Burnside lemma?


\end{document}